\unit{Why light is a quantum oscillator}
\unitref{Week 1 lectures; Week 1 tutorial}
\begin{intuition}
An optical photon is one excitation of an \emph{electromagnetic mode}.
The quantum theory does not replace Maxwell's waves with little balls:
it quantises the energy and fluctuations of each independent wave mode.
\end{intuition}
\subsection{From Maxwell's equations to a mode}
A classical electromagnetic field can be expanded in
spatial and polarisation modes. A cavity enforces standing-wave
boundary conditions; free space supports travelling-wave modes.
After suitable normalisation, the energy in one classical mode
has the harmonic-oscillator form
\[
H=\tfrac12(p^2+\omega^2q^2).
\]
Quantisation makes \(q,p\) operators satisfying
\([\hat q,\hat p]=i\hbar\). Define
\[
\hat a=\sqrt{\frac{\omega}{2\hbar}}\,\hat q
+\frac{i\hat p}{\sqrt{2\hbar\omega}},\quad
\hat a^\dagger=\sqrt{\frac{\omega}{2\hbar}}\,\hat q
-\frac{i\hat p}{\sqrt{2\hbar\omega}}.
\]
Substitution gives the central result
\[
\boxed{[\hat a,\hat a^\dagger]=1,\qquad
H=\hbar\omega\left(\hat a^\dagger\hat a+\tfrac12\right).}
\]
The zero-point term \(\hbar\omega/2\) describes the
energy left when the oscillator cannot have both conjugate
coordinates exactly zero.

\subsection{What raising and lowering actually do}
The number operator \(\hat n=\hat a^\dagger\hat a\)
has eigenstates \(\ket n\) with integer eigenvalues:
\[
\hat a\ket n=\sqrt n\ket{n-1},\qquad
\hat a^\dagger\ket n=\sqrt{n+1}\ket{n+1}.
\]
The lowest state is vacuum, \(\hat a\ket0=0\).
An arbitrary \(n\)-photon state can be built by
\[
\boxed{\ket n=\frac{(\hat a^\dagger)^n}{\sqrt{n!}}\ket0.}
\]
The factorial is required for normalisation;
it becomes especially important with two or more photons.

An observable optical electric field contains both
annihilation and creation parts:
\[
\hat E(z,t)=E_0[
\hat a e^{i(kz-\omega t)}
+\hat a^\dagger e^{-i(kz-\omega t)}].
\]
Although \(\ket n\) contains energy \(n\hbar\omega\)
above vacuum, \(\langle n|\hat a|n\rangle=0\).
It therefore has zero \emph{mean electric field}.
A nonzero mean photon number does not imply a
classically oscillating mean wave.

\subsection{Several modes and distinguishability}
For two modes \(a,b\), \(\ket{n,m}\)
means \(n\) photons in \(a\), \(m\) in \(b\).
Examples are
\[
\ket{1,1}=a^\dagger b^\dagger\ket{0,0},
\qquad
\ket{2,0}=\frac{(a^\dagger)^2}{\sqrt2}\ket{0,0}.
\]
A mode label may describe path, polarisation,
frequency, temporal envelope or a combination.
Two photons become distinguishable when a measurement
could, in principle, identify which came from which
input (even if nobody actually reads that information).
Their \emph{mode overlap} therefore controls interference.

\begin{worked}{Why \(\sqrt{2!}\) matters}
Compute
\(\|(a^\dagger)^2\ket0\|^2
=\bra0a^2(a^\dagger)^2\ket0=2\).
So \((a^\dagger)^2\ket0\) is \emph{not}
a unit vector; dividing by \(\sqrt2\)
produces \(\ket2\). This factor later changes
beam-splitter output amplitudes and probabilities.
\end{worked}
\subsection{What a phase shifter does}
Heisenberg evolution under the free oscillator Hamiltonian is
\[
\frac{d\hat a}{dt}=\frac{i}{\hbar}[H,\hat a]
=-i\omega\hat a,\qquad
\hat a(t)=e^{-i\omega t}\hat a(0).
\]
The photon number does not change: the mode
simply accumulates phase. In an interferometer,
\emph{relative} phases become measurable after
different paths are recombined.

\begin{checkpoint}
\textbf{Three different objects:} a mode specifies
which oscillator; a photon is one excitation
of that oscillator; a quantum state describes
the probabilities and coherences among its
possible excitation numbers.
\end{checkpoint}
